% problem-000047 7 丙烯醛串联反应机理
\section*{7. 丙烯醛串联反应机理}
（图：）

\subsection*{7-1}
请⽤反应机理解释产物1，3，4的⽣成过程并指出关键中间体2的构造式。

\subsection*{7-2}
化合物 5 在 Lewis 酸 $\mathrm { O _ { 3 } R e O S i P h _ { 3 } }$ 催化下苯甲醛反应得到化合物6：

（图：）

在 $\mathrm { O _ { 3 } R e O S i P h _ { 3 } }$ 催化下化合物7转化为 $\mathbf { 8 } \left( \mathrm { C _ { 1 4 } H _ { 1 8 } O _ { 2 } } \right)$ ，⽽7的异构体9则转化为10 $( \mathrm { C _ { 2 8 } H _ { 3 6 } O _ { 4 } } ) _ { \circ }$

（图：）

请画出 ${ \bf 6 } ,$ 、8和10的构造 $\updownarrow \updownarrow _ { \updownarrow }$

\subsection*{7-3}
化合物11经下列反应⽣成14：

（图：）

\subsection*{7-3}
-1 请画出化合物12，14的构造式及构型为(2�,4�,6�)-14的最稳定构象式。

\subsection*{7-3}
-2 化合物 13 如⽤ $\mathrm { T i C l } _ { 4 }$ 处理，则得到含氯的产物 15 $\left( \mathrm { C _ { 1 5 } H _ { 2 7 } C l O _ { 3 } } \right)$ ，请画出15的构造式。

\subsection*{7-4}
画出下列反应的主要产物16的构造式。

（图：）
